\section{A common price for different cuts}\label{sec:cuts}

Write $d=|V(G)|$, $e=|E(G)|$ and $w=\log q$. The variables in the
table at a vertex $v$ form a \emph{slot}. For a polynomial $f$ in
these variables, let $A(f)$ be the array of coefficients of monomials
containing exactly one variable from each slot. An index at $v$
specifies a table entry and hence a tuple of $\deg(v)$ edge labels.
We write this tuple as $\deg(v)$ separate indices of dimension $q$,
in the table basis. An isolated vertex contributes a one-dimensional
factor.

For the graph polynomial $Z_{G,q}$, a coefficient is one precisely when
the two endpoint labels agree on every edge. Thus $A(Z_{G,q})$ is a
tensor product, with one equality wire $\sum_{i=1}^q e_i\otimes e_i$
per edge. Across a fixed bipartition $C\mid V(G)\setminus C$ of the
slots its matrix rank is $q^{|\partial C|}$, where $\partial C$ is the
set of edges crossing the cut. Each crossing wire is an identity matrix
of rank $q$, and each internal wire is a nonzero vector on one side.
In the circuit and branching-program decompositions below, different
summands can have low rank across different cuts. We will use a
single tensor invariant that detects a saving on each summand and
controls the price of their sum.

\subsection*{Euclidean tensor prices}

We use the following consequences of the Euclidean price theorem
\cite[Theorem~\ext{G}{thm:main}, Section~\ext{G}{sec:spectra} and
Lemma~\ext{G}{lem:spectral-closure}]{oaG}. We work over the algebraic
closure $K=\overline F$ so that the same comparison applies to border
computations.

\begin{theorem}[Euclidean pricing]\label{thm:width}
Fix a Euclidean space of any finite dimension. Assign every index of a
finite tensor over $K$ to a site in this space, grouping all indices at
one site into a single factor. There is one compatible system of prices
$H(T)$ for all such placements. Adjoining an unused site leaves the
price unchanged. The price is real for $T\ne0$, and $H(0)=-\infty$.
\begin{enumerate}[leftmargin=*,itemsep=.3ex]
\item Tensor products add prices. Applying one linear map to the entire
factor at each site cannot increase the price.
\item Sharing a common $q$-valued label among a finite nonempty set $S$
of sites costs $\delta(S)\log q$, where
\[
       \delta(S)=\sqrt{2\pi}\,\mathbb E\max_{x\in S}\langle g,x\rangle
\]
for a standard Gaussian vector $g$. This is the Gaussian mean width
of the convex hull, normalized to give a segment its length. In particular, $\delta(\{x,y\})$
is the distance between $x$ and $y$.
\item For tensors in the same format placed at the same sites $S$,
their ordinary sum satisfies
\begin{equation}\label{eq:ordinary}
 H(T_1+\cdots+T_L)\le \max_{1\le i\le L} H(T_i)+\delta(S)\log L .
\end{equation}
\item In each fixed format and placement, the price sublevels
$\{T:H(T)\le a\}$ are Zariski closed.
\end{enumerate}
\end{theorem}

An arbitrary array on indices at a single site can be supplied by a
local linear map from a scalar. To deliver an index of dimension $q$
from that site to a site at distance $D$, contract it with one end of an
equality wire and use the other end as the output index. The wire costs
$Dw$. We can choose the auxiliary sites for this construction separately
for each summand. Only the common \emph{output} sites enter the sum
penalty in Equation~\eqref{eq:ordinary}.

For any finite nonempty set $S$ in the unit ball we have
\begin{equation}\label{eq:width-ball}
                         \delta(S)=O\bigl(\sqrt{\log(2|S|)}\bigr).
\end{equation}
Indeed each $\langle g,x\rangle$ has variance at most one. A union bound
gives $\Pr(\max_{x\in S}\langle g,x\rangle>t)\le |S|\exp(-t^2/2)$ for
$t\ge0$, and integration of the tail gives the stated estimate.

\subsection*{Fixed terminals and movable supply sites}

Give each edge its own coordinate axis in $\mathbb R^e$. Place its
two endpoint indices at the opposite unit vectors on that axis, choosing
the signs arbitrarily. Call these $2e$ sites the \emph{terminals}.
The target price at the terminals is
\[
                          H(A(Z_{G,q}))=2ew ,
\]
since each of its $e$ wires has length two. We could supply any tensor
of this format by putting its whole array at the origin and sending
each of its $2e$ indices one unit to its terminal. The graph tensor
attains this upper bound. To prove a lower bound, we will exploit a
summand's cut to supply it for strictly less.

Let $a_v$ be the sum of the terminal vectors of slot $v$. For
$C\subseteq V(G)$ set $a_C=\sum_{v\in C}a_v$. The two endpoints of
each edge internal to $C$ cancel, and the edges leaving $C$ contribute
one unit vector each on mutually orthogonal axes. Consequently
\begin{equation}\label{eq:cut}
                       \|a_C\|^2=|\partial C|.
\end{equation}
If the slots in $C$ are supplied at $z$ instead of the origin, the
distance to a unit terminal $u$ changes by at most
$-\langle z,u\rangle+\|z\|^2/2$, since
$\sqrt{1-2\langle z,u\rangle+\|z\|^2}
 \le 1-\langle z,u\rangle+\|z\|^2/2$. The change in the total
distance to the terminals belonging to $C$ is therefore at most
\begin{equation}\label{eq:move}
                -\langle z,a_C\rangle+
                    \tfrac12\|z\|^2\sum_{v\in C}\deg(v).
\end{equation}
Moving a short distance in the direction of $a_C$ saves wire length at
first-order rate $\sqrt{|\partial C|}$. The following estimate
compares that saving with the cost of sharing a rank label across the
cut, as illustrated in Figure~\ref{fig:move}.
\begin{figure}[htbp]
\centering
\begin{tikzpicture}[x=1.18cm,y=1.18cm,font=\small,
  t/.style={circle,draw=ink,fill=white,inner sep=2.4pt},
  c/.style={t,draw=sea,fill=sea!25,thick},
  wire/.style={draw=ink,thick}, cutwire/.style={draw=sea,thick},
  supply/.style={rectangle,fill=ink,inner sep=2.1pt}]
\begin{scope}[shift={(-3.45,0)}]
  \draw[wire] (-1,0)--node[midway,above]{$i$}(0,0)--node[midway,above]{$j$}(1,0);
  \node[t,label=below:$A$] at (-1,0) {};
  \node[c,label=below:$B$] at (0,0) {};
  \node[t,label=below:$D$] at (1,0) {};
  \node at (0,-1.68) {Cut $C=\{B\}$};
\end{scope}
\begin{scope}
  \draw[wire] (-1,0)--node[pos=.28,above]{$i$}(1,0);
  \draw[wire,preaction={draw=white,line width=3.8pt}] (0,-1)--node[pos=.73,right]{$j$}(0,1);
  \node[t,label=left:$A$] at (-1,0) {};
  \node[c,label=right:$B_i$] at (1,0) {};
  \node[c,label=above:$B_j$] at (0,1) {};
  \node[t,label=below:$D$] at (0,-1) {};
  \node at (0,-1.68) {Graph tensor};
\end{scope}
\begin{scope}[shift={(3.7,0)}]
  \coordinate (o) at (0,0);
  \coordinate (z) at (.5,.5);
  \draw[wire] (-1,0)--(o)--(0,-1);
  \draw[cutwire] (1,0)--(z)--(0,1);
  \draw[rust,line width=1.5pt] (o)--node[midway,below right,inner sep=1pt]{$R$}(z);
  \node[supply,label=above left:$0$] at (o) {};
  \node[supply,fill=sea,label=above right:$z$] at (z) {};
  \node[t,label=left:$A$] at (-1,0) {};
  \node[c,label=right:$B_i$] at (1,0) {};
  \node[c,label=above:$B_j$] at (0,1) {};
  \node[t,label=below:$D$] at (0,-1) {};
  \node at (0,-1.68) {Low-rank summand};
\end{scope}
\end{tikzpicture}
\caption{On the path $A$--$B$--$D$, the indices $i$ and $j$ of slot
$B$ occupy separate terminals $B_i$ and $B_j$. Each edge of the graph
gives a length-two equality wire between opposite terminals, for a
total price of $4w$. For a summand of rank at most $R$ across
$\{B\}\mid\{A,D\}$, the squares mark the two supply sites. Moving the
site $z$ toward both terminals of $B$ shortens their $q$-dimensional
wires. The wire from $0$ to $z$ has dimension $R$ and costs
$\|z\|\log R$.}\label{fig:move}
\end{figure}


\begin{lemma}[One cut per summand]\label{lem:single}
Let $G$ have $d$ vertices and $e$ edges, with
$2\le d\le w^{1/3}$, $1\le e\le 2d$ and $w=\log q$.
Fix $1\le b\le e$. Suppose $A(Z_{G,q})$ lies in the Zariski closure
of sums of at most $2^d$ tensors with the following property. For each
summand there is a bipartition of slots $C\mid V(G)\setminus C$ with
$|\partial C|\ge b$ across which that summand has matrix rank at most
$R$. Then for all sufficiently large $w$,
\[
                             \log R\ge \tfrac12 w\sqrt b .
\]
\end{lemma}

\begin{proof}
Suppose $\log R< w\sqrt b/2$, and fix one summand $T$ with its cut
$C$. A rank factorization across $C$ writes
$T=\sum_{i=1}^R U_i\otimes V_i$, with $U_i$ on the slots in $C$ and
$V_i$ on the remaining slots. Put the two supply sites at
$z=\lambda a_C/\|a_C\|$ and $0$. A single $R$-dimensional equality
wire between them supplies the common index $i$. The local maps
$e_i\mapsto U_i$ and $e_i\mapsto V_i$ produce the two sets of indices,
which are then sent to their fixed terminals. The rank wire costs
$\lambda\log R$. By Equations~\eqref{eq:cut} and~\eqref{eq:move},
the wires to the terminals save at least $\lambda w\sqrt b-ew\lambda^2$
against the original budget of $2ew$. Taking $\lambda=\sqrt b/(4e)$
gives
\[
 H(T)\le 2ew-\lambda(w\sqrt b-\log R)+ew\lambda^2
      \le 2ew-\frac{wb}{16e}.
\]

All the summands now have prices at the same $2e$ terminals. Adding
at most $2^d$ of them costs at most $O(d\sqrt{\log(2d)})$ above their
largest price by Equations~\eqref{eq:ordinary} and~\eqref{eq:width-ball}.
This is at most $wb/(32e)$ for all sufficiently large $w$, uniformly
in the given ranges, since $e\le2d$, $b\ge1$ and $d\le w^{1/3}$.
Thus the sums and their Zariski closure lie in the sublevel
$H\le2ew-wb/(32e)$. The graph tensor has price $2ew$, a contradiction.
\end{proof}
